\documentclass[11pt,a4paper]{ctexart}
\usepackage{amsmath,amssymb,geometry,booktabs,enumitem,xcolor,hyperref}
\geometry{margin=2.2cm}
\hypersetup{colorlinks=true,linkcolor=blue,urlcolor=blue}
\setlength{\parindent}{2em}
\setlist{itemsep=2pt,topsep=3pt}
\newcommand{\R}{\mathbb{R}}
\title{圆锥曲线中几何关系的统一处理方法}
\author{}
\date{}

\begin{document}
\maketitle

\section{核心思想}

把直线代入圆锥曲线方程后，总会得到一个关于参数的二次方程。若两个交点对应的参数为 $t_1,t_2$，则
\[
t_1+t_2=-\frac{\beta}{\alpha},\qquad
t_1t_2=\frac{\gamma}{\alpha}.
\]
因此，许多看似不同的问题都可以归结为三件事：
\begin{enumerate}
  \item 用参数表示曲线上的点；
  \item 把几何条件写成参数关系；
  \item 用韦达定理处理参数的和与积。
\end{enumerate}

\section{直线截圆锥曲线：一个总公式}

设圆锥曲线为
\[
F(x,y)=Ax^2+Bxy+Cy^2+Dx+Ey+F_0=0,
\]
取直线
\[
(x,y)=(x_0,y_0)+t(u,v).
\]
代入后可写成
\[
\alpha t^2+\beta t+\gamma=0.
\]
若两个交点对应 $t_1,t_2$，那么
\[
t_1+t_2=-\frac{\beta}{\alpha},\qquad
t_1t_2=\frac{\gamma}{\alpha}.
\]

当 $(u,v)$ 是单位方向向量时，$|t_1|,|t_2|$ 就是从起点到两个交点的距离。因此
\[
\boxed{|PM|\cdot|PN|=\left|\frac{\gamma}{\alpha}\right|}.
\]

其中 $\gamma=F(P)$，只由起点 $P$ 决定；$\alpha$ 只由直线方向决定。这解释了为什么长度积问题通常不必求出两个交点。

\subsection*{例：椭圆上的长度积}

对于
\[
\frac{x^2}{a^2}+\frac{y^2}{b^2}=1\qquad(a>b>0),
\]
若直线方向为单位向量 $(u,v)$，则
\[
\alpha=\frac{u^2}{a^2}+\frac{v^2}{b^2}.
\]
由于 $u^2+v^2=1$，有
\[
\frac1{a^2}\leq \alpha\leq\frac1{b^2}.
\]
所以
\[
|F(P)|b^2\leq |PM|\cdot|PN|\leq |F(P)|a^2.
\]
最大值和最小值分别对应长轴方向和短轴方向。

\section{双根、切线与极值}

当直线与圆锥曲线相切时，两个交点合并，二次方程有二重根：
\[
\boxed{\beta^2-4\alpha\gamma=0}.
\]
所以“判别式为零”在几何上就是“直线与曲线相切”。

\begin{itemize}
  \item 两个不同实根：直线是割线；
  \item 一个二重根：直线是切线；
  \item 没有实根：直线与曲线没有公共点。
\end{itemize}

长度积的最值，则来自 $t_1t_2$ 与直线方向之间的关系。

\section{用参数表示椭圆上的点}

椭圆
\[
\frac{x^2}{a^2}+\frac{y^2}{b^2}=1
\]
可参数化为
\[
P(t)=\left(
a\frac{1-t^2}{1+t^2},\quad
b\frac{2t}{1+t^2}
\right).
\]

这意味着：曲线上的点可以用一个数 $t$ 编号。复杂的几何构造就转化为参数之间的代数关系。

两点 $P(t),P(u)$ 的连线为
\[
\boxed{(1-tu)\frac{x}{a}+(t+u)\frac{y}{b}=1+tu}.
\]
特别地，它只含有 $t+u$ 和 $tu$。这正是韦达定理最适合处理的形式。

两点连线的斜率为
\[
\boxed{k_{tu}=\frac ba\cdot\frac{tu-1}{t+u}}.
\]

因此，平行、垂直、斜率比以及斜率的和与积，都可以转化为参数的有理式。

\section{复杂斜率的韦达处理}

若 $t_1,t_2$ 是某个二次方程的两个根，记
\[
s=t_1+t_2,\qquad p=t_1t_2.
\]
若某个斜率可以写成
\[
k(t)=\frac{\alpha t+\beta}{\gamma t+\delta},
\]
则
\[
k(t_1)+k(t_2)
=
\frac{2\alpha\gamma p+(\alpha\delta+\beta\gamma)s+2\beta\delta}
{\gamma^2p+\gamma\delta s+\delta^2},
\]
以及
\[
k(t_1)k(t_2)
=
\frac{\alpha^2p+\alpha\beta s+\beta^2}
{\gamma^2p+\gamma\delta s+\delta^2}.
\]
因此不需要分别求出 $t_1,t_2$。

\section{长度比：先平方}

椭圆上两点 $P(t),P(u)$ 的距离平方为
\[
|P(t)P(u)|^2
=
\frac{4(t-u)^2\left[a^2(t+u)^2+b^2(1-tu)^2\right]}
{(1+t^2)^2(1+u^2)^2}.
\]

遇到长度比时，通常先平方。例如
\[
\frac{|P(t)P(u)|}{|P(v)P(w)|}
\]
先化为两个距离平方之比，可以避免根号，并且常常只需要处理
\[
t+u,\quad tu,\quad v+w,\quad vw.
\]

\section{焦点与准线}

焦点、准线的定义为
\[
\frac{PF}{d(P,\ell)}=e,
\]
其中 $e$ 是离心率。计算时最好平方：
\[
\boxed{PF^2=e^2d(P,\ell)^2}.
\]

这样做的好处是：点到焦点的距离平方和点到直线距离平方都是二次式，代入参数后通常会得到较简单的有理式。

例如椭圆的右焦点为 $(c,0)$，右准线为 $x=a^2/c$，且
\[
e=\frac ca.
\]
对椭圆上任意点 $P$，都有
\[
PF=e\,d(P,\ell).
\]

\section{渐近线：研究无穷远方向}

设二次曲线为
\[
Ax^2+Bxy+Cy^2+Dx+Ey+F_0=0.
\]
先求出曲线的中心 $(x_0,y_0)$，令
\[
x=x_0+X,\qquad y=y_0+Y.
\]
渐近线的方向只由二次项决定：
\[
\boxed{AX^2+BXY+CY^2=0}.
\]

例如双曲线
\[
\frac{x^2}{a^2}-\frac{y^2}{b^2}=1
\]
的渐近线满足
\[
\frac{X^2}{a^2}-\frac{Y^2}{b^2}=0,
\]
所以渐近线为
\[
\boxed{y=\pm\frac ba x}.
\]

从更高的角度看，椭圆没有实的无穷远点，抛物线有一个二重无穷远点，双曲线有两个实的无穷远点，因此双曲线有两条实渐近线。

\section{调和关系：使用交比}

调和关系不是普通长度关系，而是交比关系。若 $A,B,C,D$ 在同一直线上，用有向坐标表示，则
\[
(A,B;C,D)=
\frac{c-a}{c-b}\cdot\frac{d-b}{d-a}.
\]

四点调和的定义是
\[
\boxed{(A,B;C,D)=-1}.
\]
等价地，
\[
\boxed{\frac{AC}{BC}=-\frac{AD}{BD}}.
\]
这里必须使用有向线段。

\section{极点、极线与调和关系}

设外点 $P$ 向圆锥曲线引两条切线，切点为 $A,B$。任取过 $P$ 的直线，交圆锥曲线于 $M,N$；若极线与 $MN$ 交于 $Q$，则
\[
\boxed{(M,N;P,Q)=-1}.
\]

所以，当题目同时出现切点、极线和调和关系时，应优先想到极点极线，而不是直接计算大量线段。

若用齐次坐标表示圆锥曲线
\[
\mathbf{x}^{\mathsf T}C\mathbf{x}=0,
\]
则点 $\mathbf p$ 的极线为
\[
\boxed{\mathbf{x}^{\mathsf T}C\mathbf p=0}.
\]

\section{常见几何条件的代数表达}

\begin{center}
\begin{tabular}{ll}
\toprule
几何条件 & 代数表达 \\
\midrule
三点共线 & 三点坐标行列式为零 \\
两线平行 & 斜率相等，或方向向量成比例 \\
两线垂直 & 方向向量内积为零 \\
长度相等 & 比较长度平方 \\
长度比固定 & 比较长度平方之比 \\
切线 & 二次方程有二重根 \\
焦点准线 & $PF^2=e^2d(P,\ell)^2$ \\
渐近线 & 二次项的齐次方程 \\
调和关系 & 交比等于 $-1$ \\
极线 & $\mathbf{x}^{\mathsf T}C\mathbf p=0$ \\
\bottomrule
\end{tabular}
\end{center}

\section{最后的解题流程}

遇到构造复杂的圆锥曲线问题，可以按以下顺序进行：
\begin{enumerate}
  \item 用参数 $P(t)$ 表示曲线上的点；
  \item 用直线方程或行列式表达共线、平行、共点；
  \item 长度关系先平方；
  \item 焦点准线使用距离比；
  \item 渐近线使用二次项；
  \item 调和关系使用交比；
  \item 涉及切线时使用二重根；
  \item 最后把结果化成参数的和、积或其他对称式，再用韦达定理。
\end{enumerate}

最重要的记忆是：
\[
\boxed{
\begin{array}{c}
\text{长度问题：距离平方}\\[2pt]
\text{斜率问题：参数有理式}\\[2pt]
\text{切线问题：二重根}\\[2pt]
\text{渐近线问题：无穷远方向}\\[2pt]
\text{调和问题：交比 }-1
\end{array}}
\]

这些不同形式，本质上都是在研究同一个对象：二次曲线与直线相交后产生的二次方程。

\end{document}
